% TOP_PROBLEM: {"id": 2305074, "title": "Research Problems in Function Theory — Problem 5.74", "category_id": 8, "category": {"id": 8, "name": "analysis", "display_name": "Analysis", "description": "Limits, continuity, calculus, and function theory.", "slug": "analysis", "order_index": 8, "created_at": "2026-07-31T15:26:25.671Z"}, "status": "open", "problem_number": "AMR-022-5074", "set_id": 13}
% FIRST_POSTED: 2026-10-02
% ENTRY_KIND: statement_only
% UnsolvedMath ID 2305074; source catalogue code AMR-022-5074
% !TeX program = lualatex
% Definition and sources only; this file is not an LLM solution attempt.
\documentclass[11pt,a4paper]{article}
\usepackage[margin=25mm]{geometry}
\usepackage{fontspec}
\setmainfont{lmroman10-regular.otf}[BoldFont=lmroman10-bold.otf,
 ItalicFont=lmroman10-italic.otf,BoldItalicFont=lmroman10-bolditalic.otf]
\usepackage{amsmath,amssymb,mathtools}
\usepackage{unicode-math}
\setmathfont{latinmodern-math.otf}
\AtBeginDocument{\let\setminus\smallsetminus}
\usepackage{xurl,hyperref}
\hypersetup{colorlinks=true,urlcolor=blue}
\setcounter{secnumdepth}{0}
\setlength{\parindent}{0pt}
\setlength{\parskip}{5pt}
\setlength{\emergencystretch}{3em}
\begin{document}
\section*{Research Problems in Function Theory — Problem 5.74}\label{2305074}
{\small UnsolvedMath ID 2305074 · AMR-022-5074 · Analysis · AMR Open Problem Lists}
\subsection{Definitions and mathematical statement}
Let $\mathbb T=\{\zeta\in\mathbb C:|\zeta|=1\}$ carry arclength measure and let $g:\mathbb T\to[0,\infty]$ be measurable. Characterize exactly when there is a holomorphic function $F$ on $\mathbb D$ with positive real part and with almost-everywhere boundary values $F^*$ such that
\[
g(\zeta)\le|F^*(\zeta)|\qquad\text{for almost every }\zeta\in\mathbb T.
\]
A function with positive real part has a Herglotz representation
\[
F(z)=ib+\int_{\mathbb T}\frac{\zeta+z}{\zeta-z}\,d\mu(\zeta),
\qquad b\in\mathbb R,
\]
for a finite positive measure $\mu$; a nonzero measure gives strict positivity. The boundary values in the question are finite non-tangential limits almost everywhere. Thus values of $g$ on an arclength-null set do not affect the problem.
The characterization should concern $g$ alone rather than simply reasserting the existence of $F$ or of its representing measure. Domination is by the modulus of the complex boundary value, not by its real part. Consequently the possible majorants involve both the measure and its harmonic conjugate behaviour. No bound is prescribed on $F(0)$ or on the total mass of the representing measure.
\subsection{Short English statement}
Which nonnegative boundary functions are almost everywhere dominated by the modulus of a disc analytic function with positive real part?
\subsection{Sources}
\begin{itemize}
\item[S1] ULAM.AI and the UnsolvedMath Contributors, supplied problems.json, record 2305074 (AMR-022-5074), AMR Open Problem Lists. Catalogue identity and source transcription; CC BY 4.0.
\url{https://huggingface.co/datasets/ulamai/UnsolvedMath}.
\item[S2] Hayman - Research Problems in Function Theory (2018), Problem 5.74. Original source indexed by AMR-022-5074.
\url{https://arxiv.org/abs/1809.07200}.
\end{itemize}
\end{document}
